\chapter{A frozen detector-calibration experiment}
\label{ch:dc-experiment}

The finite certificate was evaluated on $1{,}600$ independently seeded
synthetic datasets, with unknown detector gain, exposure width and diffusion
scale carried through the inference.  The driven case produced a positive
geometry certificate in all $400$ repetitions; equilibrium produced none.
The acquisition is expensive: each dataset represents more than $80$ million
detector readings.  Higher polynomial order did not improve the selected
bound in this design.  These results demonstrate a working finite inference
chain under its declared assumptions, with its sampling and calibration
costs retained.  They are not laboratory observations or an efficiency claim.

\section{The frozen acquisition and inference}

The data-producing process was the stationary circle diffusion
$d\phi_t=\omega\,dt+\sqrt{2\gamma}\,dW_t$, observed through
$f=\cos\phi$.  Three uniform exposure windows start at $0,h,2h$ and have
width $\tau$.  Each binary detector output $R\in\{-1,1\}$ has conditional
mean $\kappa A$, where $A$ is the average of $f$ over its window.  The two
endpoint windows have one readout each; the center has four conditionally
independent readouts of the same exposure.  Translating all window times
makes this the causal-boxcar convention of the preceding theorem, without
changing its stationary moment law.
The detector is noninvasive in this model.  No noise variance is supplied
to inference; bounded outputs and the conditional replica contract justify
the moment calculation.

The design was fixed at
\[
 h=7/200=0.035,\quad \lambda=1,\quad m_{\max}=1,
 \quad\kappa_0=0.9,\quad n=13{,}333{,}333.
\]
Every dataset also has $100{,}000$ constant-signal references and $100{,}000$
step references.  The constant reference has mean $\kappa$; the step is read
$s=1/500{,}000$ after its known transition and has mean $\kappa s/\tau$.
Inference starts from $\kappa\in[0,1]$ and
$\tau\in[1/200{,}000,1/50{,}000]$.  It receives the common physical bounds
$M_{1,+}=M_{1,-}=4$ and $M_{2,+}=M_{2,-}=13$.

The gain and width supplied to the simulator are not supplied as exact
inference inputs.  The diffusion scale is likewise unknown to the
certificate: its upper bound comes from the covered twin and lag moments.
The pointwise identity $\Gamma_f=\gamma(1-f^2)$ and the generator caps remain
assumptions.  The inference does not fit the complete circle model or
substitute its arcsine marginal as a measured second moment.

\begin{table}[htbp]
\centering
\begin{tabular}{lrrrrr}
\hline
Case & $\gamma$ & $\omega$ & $\kappa$ & $10^5\tau$ & $\sigma$\\
\hline
Driven & 1 & 3 & 0.9 & 1 & 9\\
Equilibrium & 1 & 0 & 0.9 & 1 & 0\\
Weak drive & 1 & 1 & 0.9 & 1 & 1\\
Changed calibration & 0.5 & 3 & 0.8 & 1.5 & 18\\
\hline
\end{tabular}
\caption{Simulator parameters and exact entropy $\sigma=\omega^2/\gamma$.
Each case has $400$ datasets.  The changed-calibration case uses fresh
references with altered gain, exposure and diffusion; it is not an
uncalibrated change of detector between target and reference.}
\label{tab:dc-benchmark-cases}
\end{table}

Two-sided empirical-Bernstein intervals cover the fixed target-statistic
vector at total failure allocation $0.04$.  Bounded reference-mean intervals
receive allocation $0.01$.  The resulting $0.05$ per-dataset allowance covers
gain and shutter projection, the geometry scale, and selection among the
predeclared order-zero and order-one bounds.  The fixed nominal-gain
contrasts retain covariance before calibration projection; componentwise
bounds are also retained, and their covered maximum is reported.

The endpoint bound~\eqref{eq:dc-endpoint-baseline} uses the same records,
reference intervals and simultaneous event.  Here $b=1-\lambda h=0.965$.
It is an established reversible-support and entropy-bridge comparison.  It
does not require the pointwise geometry identity, so it has broader
structural scope.  Comparing the two values measures what the additional
geometry assumption provides in these examples; it does not establish
dominance under identical assumptions or against all inference methods.

\section{A finite-law simulator with an explicit error charge}

The retained six-bit law is a finite exposure-moment polynomial:
\begin{equation}
 \Pr(R=r)=2^{-6}\mathbb E\prod_{i=1}^{6}
                (1+\kappa r_i A_{w(i)}).
 \label{eq:dc-bit-law}
\end{equation}
The simulator computes the needed integrated Fourier moments, rather than
approximating a Brownian path on a time grid.  A closed system with $140$
real states supplies the moments through total degree six.  Rational Taylor
actions enclose the exposure and free-gap semigroups; if
$x=t\|G\|_\infty$, the degree-$N$ action remainder is bounded by
\[
 \|v\|_\infty\,3^{\lceil x\rceil}\frac{x^{N+1}}{(N+1)!}.
\]
Input uncertainty, exposure normalization and event-polynomial coefficients
are propagated explicitly.  In the driven case the exposure stages use
Taylor degree $15$ and the gaps degree $58$.

The resulting law is rounded to a nonnegative $63$-bit dyadic probability
vector summing exactly to one.  Conditional center exchangeability collapses
the $64$ labeled outcomes to $20$ categories: left sign, center plus-count,
and right sign.  These counts retain all predeclared symmetric statistics;
they would not support an arbitrary later score that distinguishes labeled
center channels.  Integer-CDF sampling avoids floating probability thresholds.
Independent uniform random bits are the ideal sampling assumption; the
implementation uses NumPy $2.3.5$ and PCG64 pseudorandom streams.

The exact rational total-variation charges include Taylor enclosure,
dyadic rounding, all target records and both reference batches.  Rounded
upward decimal bounds per dataset are
\[
 1.308\,10^{-11},\quad 1.717\,10^{-11},\quad
 1.067\,10^{-11},\quad 1.141\,10^{-11}
\]
in the order of Table~\ref{tab:dc-benchmark-cases}.  The full bank's product-law
charge is below $2.093\,10^{-8}$.  Thus transferring a per-dataset probability
statement between the ideal model and the sampled dyadic law adds its
specified charge.  These bounds concern the probability-law approximation;
they do not certify hardware or the literal randomness of a deterministic
pseudorandom generator.

\section{Results, inclusion checks and negative controls}

\begin{table}[htbp]
\centering
\small
\begin{tabular}{lrrrr}
\hline
Case & Geometry positive & Mean lower & Endpoint positive & Mean lower\\
\hline
Driven & $400/400$ & $0.57735$ & $400/400$ & $8.1640\,10^{-5}$\\
Equilibrium & $0/400$ & $0$ & $0/400$ & $0$\\
Weak drive & $0/400$ & $0$ & $27/400$ & $6.8211\,10^{-9}$\\
Changed calibration & $395/400$ & $0.59905$ & $400/400$ & $6.9611\,10^{-5}$\\
\hline
\end{tabular}
\caption{Observed synthetic detection counts and mean lower bounds.  The
geometry and endpoint bounds use the same detector records but different
structural assumptions.  No method exceeded the true entropy in these runs.}
\label{tab:dc-benchmark-outcomes}
\end{table}

The geometry lower bounds in the driven case have median $0.57649$ and
empirical $2.5$th--$97.5$th percentiles $[0.29974,0.84860]$.  Under changed
calibration the median is $0.59490$; the corresponding percentiles are
$0.05523$ and $1.14799$.  These are descriptive quantiles across repeated
datasets, not confidence intervals for the true entropy.  The mean
certificates represent only about $6.4\%$ and $3.3\%$ of the respective
entropies $9$ and $18$, despite the large acquisition quota.

\begin{table}[htbp]
\centering
\small
\begin{tabular}{lrrrr}
\hline
Case & Target event & Gain/width inclusion & Finite $\gamma$ bounds & Mean upper\\
\hline
Driven & $400/400$ & $400/400$ & $400/400$ & $1.69770$\\
Equilibrium & $399/400$ & $400/400$ & $400/400$ & $1.53917$\\
Weak drive & $399/400$ & $400/400$ & $400/400$ & $1.55591$\\
Changed calibration & $397/400$ & $400/400$ & $400/400$ & $1.20801$\\
\hline
\end{tabular}
\caption{Independent replay of the simultaneous target-mean event,
reference-region inclusion and geometric-scale output.  Every finite
$\gamma$ upper bound covered its simulator value; there were no unbounded
geometry abstentions.  The target event missed five times in total without
causing an entropy overshoot.}
\label{tab:dc-benchmark-coverage}
\end{table}

All $1{,}600$ count records were independently reconstructed and replayed.
The target mean-event inclusion check used rational law enclosures and had
no unresolved enclosure-boundary cases.  Calibration inclusion and entropy
overshoot checks are separate: failure of a sufficient simultaneous event
need not cause the final lower bound to overshoot.

The equilibrium control gives no positive result, as required for a useful
check against obvious false positives.  The weak-drive case is an equally
important negative result for sensitivity: the geometry certificate is zero
throughout, while the endpoint certificate is positive in $27$ runs.  Under
changed calibration the endpoint method also returns five more positive
results.  The geometry method therefore does not dominate even this one
broader-class comparator in detection frequency.

A transparent derived fallback is to report the maximum of the geometry and
endpoint lower bounds.  Under the narrower geometry contract, both hold on
the already covered event, so this maximum requires no additional error
allocation.  Its positive counts are $400,0,27,400$ in the four cases.
This is a deterministic combination made after the frozen benchmark, not a
change to its per-method outputs or a separately preregistered experiment.

Zero overshoots in $400$ independent repetitions do not prove zero failure
probability.  Under the ideal repeated-sampling model, their one-sided
$95\%$ binomial upper bound is $1-0.05^{1/400}\simeq0.00746$, separately for
each case and method.  This diagnostic is not a simultaneous guarantee
across the comparisons.  The finite-sample coverage claim comes from the
preceding theorem and implementation, not from the observed frequencies.

\section{What the higher-order and design comparisons did not show}

The selected prefix was order zero in every dataset.  Order one improved
neither the final driven result nor the changed-calibration result.  Retaining
covariance did help its own unselected estimate: in the driven case the
order-one mean rose from $0.02109$ for the componentwise calculation to
$0.06497$ for the fixed-contrast projection.  Both remained below the selected
order-zero bound.  The sharp infinite hierarchy is a theorem; useful
finite-resource gains from its higher orders were not demonstrated here.

Before freezing the experiment, deterministic population calculations
examined lags and replica counts.  At $h=0.035$, under a ceiling of $80$
million target readouts, two, four and six center replicas permit respectively
$20{,}000{,}000$, $13{,}333{,}333$ and $10{,}000{,}000$ stationary records.
The corresponding population-moment values inserted into the confidence
formula were $0.54161$, $0.55834$ and $0.50535$, with order zero selected in
all three.  All designs add $200{,}000$ reference readouts.  The prospective
calculation used target/reference allocations $0.025/0.025$, distinct from
the final $0.04/0.01$ allocation.

Those numbers are design diagnostics, not confidence-certified outcomes,
power estimates or additional sampled experiments.  They show why comparing
replicas at fixed record count would conceal extra readout costs.  Earlier
$h=0.02$ diagnostics returned zero for the strong-drive geometry score at
two and ten million target records, and higher orders remained zero even at
fifty million.  A separately retained design with narrower exposure and
tighter supplied generator caps looked better, but changed the apparatus
contract and was not used for the primary experiment.  Its gains cannot be
attributed to the estimator alone.

\begin{figure}[htbp]
\centering
\includegraphics[width=\linewidth]{figures/detector_calibration.pdf}
\caption{Left: distributions over the $400$ simulated datasets in each case,
with boxes at the interquartile range and whiskers at empirical $2.5$th and
$97.5$th percentiles; the vertical axis uses a symmetric logarithmic scale.
Black diamonds mark the simulator entropy.
The endpoint bound has broader structural validity than the geometry bound.
Right: the separate prospective population diagnostic at matched target
readout ceilings, with the additional reference readings counted.  This
panel contains no random experimental outcomes or confidence certificates.}
\label{fig:dc-benchmark}
\end{figure}

\section{Acquisition costs and reproducibility}

Each dataset uses $79{,}999{,}998$ target readings, $200{,}000$ reference
readings and $80{,}199{,}998$ readings in total.  It requires
$13{,}333{,}333$ independent stationary preparations and $200{,}000$ controlled
reference preparations.  Four center shots share one physical exposure;
they still count as four detector readings.  There are $39{,}999{,}999$
target windows and $200{,}000$ reference windows.

The summed modeled target evolution duration is $n(2h+\tau)$:
$933{,}466.64333$ model-time units for the first three cases and
$933{,}533.309995$ for changed calibration.  The summed target exposure
durations are respectively $399.99999$ and $599.999985$; reference exposure
durations add $2$ and $3$.  These are sums over independent acquisitions,
not wall-clock durations or measured energy.  Stationary preparation, reset,
replica construction, clock calibration and shutter-control costs remain
unmeasured.  Fast categorical simulation does not make those physical
preparations inexpensive.

The authoritative files in \texttt{experiment/} are \texttt{contract.json},
\texttt{benchmark.json}, the compressed sufficient counts
\texttt{benchmark\_counts.json.gz}, and the first-generation
clock record \texttt{benchmark\_run.json}.  The contract binds
simulator, certifier and benchmark source hashes.  It was frozen on
2026-10-13 at 01:45:27 UTC; generation started at 01:45:40 and completed at
01:48:01 with four workers.  The choices preceding this freeze are retained
as prospective population diagnostics, not concealed as a preselected
universal design.  No stochastic pilot target records were used.

The seed rule is
$1032026400000+100000\,i+10\,r+s$ for case index $i$, repetition $r$ and
stream $s\in\{0,1,2\}$ for target, gain and step records.  Exact counts,
law enclosures, source bindings and the independent replay provide the
reproduction chain.  The stochastic-drive family proved in the preceding
chapter was not simulated in this bank.  Its analytic result and this
constant-drift detector benchmark remain distinct evidence.
